\chapter{Conclusiones y trabajo futuro}
\label{chap:conclusiones}

Los resultados obtenidos en la experimentación permiten cerrar el objetivo de este trabajo, contrastar los hallazgos empíricos con los objetivos formulados en la introducción, valorar qué implican para el estado actual del campo y señalar qué queda por explorar. Este capítulo recoge, en su primera sección, las conclusiones derivadas de la comparación entre Word2Vec y QWord2Vec. La segunda sección propone las posibles líneas de investigación a partir de lo realizado en este trabajo.

\section{Conclusiones}

Este Trabajo de Fin de Grado se ha planteado con un objetivo doble como se indicó en la introducción, implementar y evaluar empíricamente el modelo cuántico QWord2Vec frente a su contraparte clásica Word2Vec, y, sobre esa base, responder a una pregunta más amplia: si el espacio de Hilbert puede actuar como espacio latente para los \textit{embeddings} de palabras con resultados equiparables a los del espacio euclídeo clásico. La aproximación seguida se separa de la del trabajo de referencia \cite{ohno_qword2vec_2025} en un punto esencial. Mientras dicho estudio entrenaba ambos modelos bajo condiciones idénticas para medir hasta qué punto la versión cuántica reproducía las representaciones de la clásica, aquí se ha optado por aplicar a cada modelo las mejores técnicas disponibles dentro de su propio paradigma, con el fin de evaluar qué puede ofrecer cada uno cuando se le permite operar en su mejor versión.

Sobre esta premisa metodológica se han cumplido los objetivos formulados en el Capítulo \ref{chap:introduccion}.

El primer hallazgo relevante atañe a la calidad semántica de las representaciones aprendidas. Word2Vec ha obtenido un \textit{cosine\_delta} de $0.8927$ en la fase inicial y de $0.9360$ en la fase ampliada, evidenciando una clara separación entre los clústeres semánticos definidos sobre el corpus. QWord2Vec, por su parte, ha alcanzado un \textit{cosine\_delta} de $0.0550$ en la fase inicial y de $0.3293$ en la fase ampliada. Al emplear la misma métrica en ambos modelos, la comparación es directa, existe una diferencia significativa en favor de Word2Vec, si bien la brecha se reduce de forma apreciable al pasar al corpus equilibrado. Este contraste cobra mayor relevancia al considerar que QWord2Vec opera con menos información de entrenamiento: mientras Word2Vec aprende a partir de todos los pares (palabra, contexto) extraídos del corpus, el modelo cuántico se entrena únicamente con los dos contextos más frecuentes de cada palabra objetivo. La mejora sustancial del \textit{cosine\_delta} entre fases sugiere que la calidad de los \textit{embeddings} cuánticos es especialmente sensible al equilibrio del corpus, y apunta a que la brecha podría reducirse con condiciones de entrenamiento más favorables, lo que mantiene el interés del espacio de Hilbert como espacio latente para representaciones lingüísticas en línea con la pregunta central de este trabajo.

A esta ventaja se le suma una segunda ventaja, la cual se ve proyectada principalmente al escalar el problema. Word2Vec almacena las representaciones en dos matrices de tamaño $V \times d$ (con $V$ el tamaño del vocabulario y $d$ la dimensión del \textit{embedding}), de modo que el coste en parámetros crece de forma lineal con el número de palabras. QWord2Vec, en cambio, codifica cada palabra como un estado cuántico sobre $n$ qubits que admite hasta $2^n$ estados distintos en superposición, y los parámetros entrenables del circuito escalan como $\mathcal{O}(L \cdot n)$, es decir, de forma logarítmica en el tamaño del vocabulario. En el corpus utilizado en este trabajo, con apenas $13$ palabras, esta ventaja no llega a materializarse, la diferencia absoluta de parámetros es pequeña y queda oculta por el coste de la simulación clásica. Sin embargo, a medida que crece $V$, la representación cuántica requeriría progresivamente menos memoria que la clásica para almacenar y manipular los \textit{embeddings}, propiedad que constituye uno de los argumentos teóricos más sólidos a favor del paradigma cuántico para tareas de \gls{pln} a gran escala.

La correlación de Pearson entre los vectores de distancias coseno de ambos modelos completa esta lectura, pasando de $0.2920$ en la fase inicial a $0.4940$ con el corpus equilibrado. Ambos valores quedan por debajo del $0.81$ reportado en el trabajo de referencia, pero esa distancia no debe interpretarse como un retroceso, sino como una consecuencia directa del objetivo planteado. El estudio original buscaba precisamente que las dos representaciones fueran lo más parecida posible, mientras que aquí cada modelo se ha entrenado para optimizar lo mejor posible su propio espacio. Es por esto que las geometrías resultantes en el espacio euclídeo y en el de Hilbert no tienen por qué coincidir. La mejora observada entre fases sugiere, además, que un corpus más rico y equilibrado favorece de forma natural una mayor alineación entre ambas representaciones.

Junto a la calidad de los \textit{embeddings}, los resultados obtenidos confirman las propiedades teóricas atribuidas al optimizador \gls{qn-spsa}. El modelo cuántico ha convergido en aproximadamente $360$ iteraciones, frente a las hasta $20\,000$ empleadas en el trabajo de referencia, gracias a la combinación de un ajuste fino de los hiperparámetros mediante \textit{Grid Search} en dos fases, una inicialización por \textit{educated guess} sobre diez candidatos aleatorios y la activación del parámetro \textit{blocking}. Junto a esta convergencia acelerada se ha observado que el componente de \gls{qng} permanece inactivo cuando el problema es de pequeña escala. El modelo alcanza la convergencia antes incluso de que dicho tensor comience a aplicarse, de modo que no llega a intervenir en las iteraciones efectivas y el optimizador se comporta en la práctica como un \gls{spsa} bien ajustado. Por tanto, los resultados obtenidos no permiten valorar su contribución real, sino únicamente constatar que en este problema en específico no llega a ejecutarse. Cabe esperar que su aporte sea apreciable en problemas con vocabularios extensos y circuitos más profundos, donde la convergencia requiera un número de iteraciones suficiente para que el tensor entre en juego.

Un resultado llamativo observado durante la experimentación es la extrema sensibilidad de QWord2Vec a la semilla aleatoria del simulador cuántico. Cambiar la semilla no produce variaciones menores en el resultado final, sino que la función de pérdida, el \textit{error rate}, la velocidad de convergencia y la calidad semántica de los \textit{embeddings} pueden diferir de forma tan pronunciada entre ejecuciones que el comportamiento equivale prácticamente al de un modelo distinto en cada caso. Identificar la semilla que da lugar a un buen entrenamiento requiere, por tanto, explorar varias ejecuciones, algo que contrasta con la estabilidad reproducible de Word2Vec. La raíz de este fenómeno se encuentra en la acumulación de dos fuentes independientes de aleatoriedad que se potencian entre sí. Por un lado, la estimación de las distribuciones de probabilidad del circuito se realiza mediante un número finito de \textit{shots} (2048) por circuito, los valores de probabilidad medidos son estimaciones estadísticas ruidosas de los valores ideales, y ese ruido se propaga directamente a la función de pérdida. Por otro lado, \gls{qn-spsa} estima el gradiente perturbando aleatoriamente los parámetros del circuito en cada iteración, de modo que la dirección de actualización depende a su vez de la semilla. La interacción entre el ruido de muestreo cuántico y las perturbaciones estocásticas del optimizador genera trayectorias de optimización que pueden divergir drásticamente desde las primeras iteraciones, llevando a mínimos locales de calidad muy diferente. Este fenómeno es inherente a los algoritmos variacionales cuánticos en regímenes de pocos \textit{shots} y constituye uno de los principales retos abiertos para su uso en aplicaciones reproducibles.

Estos resultados deben situarse, no obstante, dentro de las limitaciones que los condicionan. La más determinante es la simulación clásica de los circuitos cuánticos, cuyo coste crece exponencialmente con el número de qubits e impone un techo práctico al tamaño del corpus manejable en un equipo personal. Su impacto se traduce de forma directa en los tiempos de ejecución medidos: aproximadamente $1$ segundo para Word2Vec frente a unos $50$ minutos para QWord2Vec, una diferencia de cuatro órdenes de magnitud que no responde a la complejidad algorítmica del modelo, sino al carácter clásico del simulador. Por último, conviene recordar la diferencia de madurez entre ambos paradigmas, Word2Vec lleva más de una década consolidándose como referencia del \gls{pln}, mientras que QWord2Vec se publicó hace menos de un año, por lo que el cuerpo de optimizaciones, librerías especializadas y conocimiento práctico disponibles para el modelo cuántico es todavía muy limitado.

Considerando los resultados y las limitaciones de manera conjunta, QWord2Vec demuestra aprender representaciones semánticamente coherentes, aunque su \textit{cosine\_delta} ($0.0550$ en la primera fase y $0.3293$ en la segunda) queda por debajo del obtenido por Word2Vec ($0.8927$ y $0.9360$ respectivamente). La mejora observada entre fases al equilibrar el corpus es prometedora, pero no alcanza todavía la calidad semántica ni la eficiencia computacional del modelo clásico. Antes que una valoración definitiva sobre el potencial de la computación cuántica en \gls{pln}, este Trabajo de Fin de Grado debe entenderse como un análisis del estado actual del campo: una tecnología con fundamentos teóricos sólidos y resultados iniciales que despiertan interés, pero que requiere todavía avances importantes, tanto en hardware como en algoritmia, para competir con los estándares clásicos en condiciones realistas.


\section{Trabajo a futuro}

Las conclusiones obtenidas dejan abiertas diversas líneas de investigación que prolongan de forma natural lo realizado en este trabajo. Estas líneas pueden agruparse en tres grandes ejes: la escalabilidad del modelo, la aplicación de los \textit{embeddings} a tareas posteriores y la exploración de variantes algorítmicas y arquitectónicas.

El primer eje, y probablemente el más urgente y evidente, es la validación de QWord2Vec sobre corpus de mayor escala. Los experimentos descritos se han limitado a vocabularios de $13$ palabras debido al coste exponencial de la simulación clásica de circuitos cuánticos y que no se tiene la capacidad de cómputo local para escalar más el problema. Es decir, para superar esta barrera será necesario recurrir a simuladores desplegados en equipos de computación de alto rendimiento, en concreto a hardware cuántico real, particularmente a dispositivos \gls{nisq}. Esta validación permitiría además observar si las ventajas teóricas que se conocen del espacio de Hilbert se traducen efectivamente en mejoras prácticas cuando el problema crece, comparándolo con los resultados obtenidos en simulación.

El segundo eje, tal y como se señala en el artículo de referencia seguido \cite{ohno_qword2vec_2025}, es la aplicación de los \textit{embeddings} cuánticos a tareas posteriores. La calidad de una representación vectorial no se mide únicamente por sus propiedades intrínsecas como se ha realizado en este TFG, sino que se lleva más allá midiéndolo sobre todo por su utilidad en tareas como clasificación de texto, agrupamiento, detección de similitud o recuperación de información. Resulta particularmente atractivo el desarrollo de modelos generativos cuánticos que aprovechen estas representaciones como entrada. Esta integración con tareas reales permitiría además identificar qué propiedades del espacio de Hilbert resultan verdaderamente diferenciales en escenarios aplicados.

El tercer eje agrupa las mejoras posibles sobre el optimizador y la arquitectura del circuito. En primer lugar, convendría comprobar si el tensor métrico de Fubini-Study aporta valor en problemas de mayor complejidad En segundo lugar, la alta sensibilidad a la semilla aleatoria documentada en este trabajo abre la puerta a estudiar estrategias de entrenamiento más robustas, como ejecutar múltiples inicializaciones con distintas semillas y seleccionar la mejor solución, aumentar el número de \textit{shots} para reducir el ruido de muestreo, o incorporar técnicas de promediado de trayectorias que amortigüen la varianza entre ejecuciones. En tercer lugar, existe una técnica de medición parcial sobre qubits \cite{huang_power_2021}, también mencionada en el artículo de referencia \cite{ohno_qword2vec_2025}, que permitiría aumentar la no linealidad del modelo y reducir la probabilidad de error, mejorando potencialmente la calidad de los \textit{embeddings} sin necesidad de incrementar el número de qubits del sistema. Por último, la exploración de arquitecturas alternativas del circuito parametrizado, variando la profundidad, la disposición de las puertas de rotación o la estrategia de entrelazamiento, con el objetivo de encontrar configuraciones que ofrezcan un mejor resultado de expresividad y entrenabilidad.

Con todo esto se dibuja una hoja de ruta coherente para que QWord2Vec, el primer circuito cuántico propuesto específicamente para este problema, pueda dar el salto desde su validación inicial sobre corpus reducidos hasta una posible aplicación realista en \gls{pln}. Avanzar en ellas requerirá tanto el progreso del hardware cuántico como el desarrollo de nuevas técnicas algorítmicas adaptadas a este tipo de modelos. Solo cuando ambos frentes maduren será posible determinar si las ventajas teóricas que la computación cuántica ofrece para el procesamiento del lenguaje natural se traducen en mejoras tangibles y reproducibles sobre los estándares clásicos.
